\section{Introduction}\label{sec:intro}

A model that ranks stocks by next-period return is judged by two numbers: its information coefficient (IC), the cross-sectional correlation between scores and realised returns, and the net Sharpe ratio of a book built from the scores. The first is well documented and mostly encouraging \citep{gkx2020,chenpelgerzhu2024,kellymalamudzhou2024}; the second is less so \citep{novymarxvelikov2016,avramov2023}. The fundamental law of active management links them \citep{grinold1989,clarke2002}: a Sharpe ratio is an IC times a transfer coefficient set by the portfolio construction, less a cost drag set by turnover. We measure each term on one dataset under a pre-registered protocol and answer an operational question: can a model built from free data reach a net Sharpe ratio of 1.5 with drawdown under 25\%, and do graph methods help?

The data are deliberately modest: a public survivor-only price feed, SEC XBRL fundamentals, 8-K earnings dates, SIC and BEA input--output industry graphs, customer--supplier links from filing text; nothing was purchased. The study is therefore reproducible by a student and carries a survivorship bias that we bound but cannot remove (Section~\ref{sec:limits}).

\paragraph{Contributions.}
\begin{enumerate}
\item A pre-fixed shallow booster carries a rank signal whose sign replicates out of its development window and on one held-out year; per-refit hyperparameter selection on a noisy validation slice removes its advantage over a linear control.
\item Along a 55-cell grid of constructions, an exact cost-ladder identity shows that the shortfall against the target is information, not cost: no gross Sharpe ratio reaches 1.5, so no cost level closes the gap, and the least demanding cell would need, to first order, at least about $\numKappaLeastBoth\times$ its IC.
\item Relational graph networks trained end to end on four point-in-time firm graphs add a paired $\Delta\IC$ against an identical no-graph network that is indistinguishable from zero; the rewired-graph comparison, completed only as monthly and quarterly IC levels, points the same way \citep{maslovsneppen2002}.
\item A power calculation places confirmation of a Sharpe-0.3 property decades away, so the sizing decision rests on a prior.
\end{enumerate}

\paragraph{Related work.}
Trees and shallow networks are the strongest cross-sectional predictors on CRSP, with small out-of-sample $R^2$ and large pre-cost decile spreads \citep{gkx2020}, much of which costs and microcaps remove \citep{avramov2023}. Economic links predict firm returns \citep{cohenfrazzini2008,menzlyozbas2010}; graph neural networks \citep{kipfwelling2017,hamilton2017,velickovic2018,schlichtkrull2018} are their end-to-end form; we know of few tests of them against a rewired-graph null on a point-in-time panel, and ours is partial (Section~\ref{sec:limits}). Inference under selection follows \citet{harveyliuzhu2016}, \citet{baileylopez2014} and \citet{hansen2005}, Sharpe-ratio uncertainty \citet{lo2002} and \citet{ledoitwolf2008}, and the decay of published anomalies \citep{mcleanpontiff2016,houxuezhang2020} is the prior for reading a small signal.

Figure~\ref{fig:pipeline} (p.~\pageref{fig:pipeline}) gives the pipeline in one picture. Sections~\ref{sec:data}--\ref{sec:infer} define the panel, validation geometry, predictors, graphs, book construction and statistics; Sections~\ref{sec:results}--\ref{sec:limits} report results, power and limitations.
